\documentclass[a4paper, serif, 10pt]{moderncv}

%% moderncv
% color and style
\moderncvcolor{black}
\moderncvstyle{banking}

%% page layout/margins
\usepackage[top=2.5cm, bottom=2.5cm, left=1.5cm, right=1.5cm]{geometry}

\nopagenumbers{}

%% Babel config for German language
% ref:
% - https://latex3.github.io/babel/guides/locale-german.html
\usepackage[ngerman]{babel}

%% fontspec
\usepackage{fontspec}

\IfFontExistsTF{Linux Libertine}{
  \setmainfont[Ligatures={TeX, Common}, Numbers=OldStyle]{Linux Libertine}
}{}
\IfFontExistsTF{Linux Libertine O}{
  \setmainfont[Ligatures={TeX, Common}, Numbers=OldStyle]{Linux Libertine O}
}{}

%% CTeX
% CJK support, used to show CJK characters in the resume
%
% - heading=false: disable ctex heading
% - punct=kaiming: use kaiming punctuations styles for CJK
% - scheme=plain: use plain scheme, do not override \normalsize font size
% - space=auto: space settings for CJK characters
%
% CJK fonts are resolved per locale language, e.g. Noto Serif/Sans CJK JP
% for Japanese resumes, with fallback to the next region if not installed
%
% ref:
% - http://ctan.mirrorcatalogs.com/language/chinese/ctex/ctex.pdf
\usepackage[UTF8, heading=false, punct=kaiming, scheme=plain, space=auto]{ctex}

\IfFontExistsTF{Noto Serif CJK SC}{
  \setCJKmainfont{Noto Serif CJK SC}
}{}
\IfFontExistsTF{Noto Sans CJK SC}{
  \setCJKsansfont{Noto Sans CJK SC}
}{}

%% line spacing
\usepackage{setspace}
\setstretch{1.125}

%% URL styling - use normal text instead of monospace
\urlstyle{same}

%% Patch \httplink and \httpslink to handle full URLs
% The original moderncv macros always prepend http:// or https:// to URLs.
% This causes issues when the URL already has a protocol (e.g., https://example.com)
% resulting in malformed URLs like https://https://example.com.
% This patch checks if the URL already contains "://" and uses it directly if so.
% Note: We use \str_set:Nx (x-expansion) to expand macros like \@homepage first.
\ExplSyntaxOn
\str_new:N \g_yamlresume_url_str

\RenewDocumentCommand{\httpslink}{O{}m}{
  \str_set:Nx \g_yamlresume_url_str {#2}
  \str_if_in:NnTF \g_yamlresume_url_str {://}
    { \url{#2} }
    { \url{https://#2} }
}

\RenewDocumentCommand{\httplink}{O{}m}{
  \str_set:Nx \g_yamlresume_url_str {#2}
  \str_if_in:NnTF \g_yamlresume_url_str {://}
    { \url{#2} }
    { \url{http://#2} }
}
\ExplSyntaxOff

\name{Dr. Lena Hartmann}{}
\title{Senior KI-Ingenieur – Maschinelles Lernen \& Generative KI}
\phone[mobile]{+49 30 12345678}
\email{lena.hartmann@example.com}
\homepage{https://lena-hartmann.ai}

\address{Friedrichstraße 123, Berlin, Berlin, Deutschland, 10117}{}{}

\extrainfo{{\small \faGithub}\ \href{https://github.com/lena-hartmann}{@lena-hartmann} {} {} {} • {} {} {} 
{\small \faLinkedin}\ \href{https://www.linkedin.com/in/lena-hartmann}{@lena-hartmann} {} {} {} • {} {} {} 
{\small \faStackOverflow}\ \href{https://stackoverflow.com/users/1234567/lena-hartmann}{@lena-hartmann}}

\begin{document}

\maketitle

\section{Basisinformationen}

\cvline{}{\begin{itemize}
\item Senior KI-Ingenieur mit über 8 Jahren Erfahrung in der Entwicklung und produktiven Skalierung von Systemen für Maschinelles Lernen und Generative KI.
\item Spezialisiert auf Deep Learning, Sprachverarbeitung und MLOps; verantwortlich für End-to-End-Pipelines von der Datenaufbereitung bis zum Monitoring im Betrieb.
\item Nachweisliche Erfolge bei der Führung interdisziplinärer Teams und der Übersetzung komplexer Geschäftsanforderungen in robuste, ethisch verantwortungsvolle KI-Lösungen.
\item Leidenschaft für innovative Forschungsansätze, reproduzierbare Experimente und die kontinuierliche Verbesserung von Modellqualität und -performance.
\end{itemize}}
\section{Ausbildung}

\cventry{Okt. 2012–Sept. 2016}
        {Doktor, Künstliche Intelligenz und Maschinelles Lernen, Punkte: summa cum laude}
        {Technische Universität Berlin}
        {\url{https://www.tu.berlin}}
        {}
        {\begin{itemize}
\item Promotion über selbstüberwachtes Lernen für die Sprachverarbeitung mit Fokus auf ressourcenarme Domänen.
\item Entwicklung neuartiger Architekturen, die auf mehreren Benchmarks state-of-the-art Ergebnisse erzielten.
\item Publikationen auf führenden Konferenzen wie NeurIPS, ICML und ACL.
\end{itemize}
\textbf{Kurse}: Deep Learning, Probabilistische Graphenmodelle, Natural Language Processing, Reinforcement Learning, Bayesianische Inferenz}

\cventry{Okt. 2010–Sept. 2012}
        {Master, Informatik, Punkte: 1,2}
        {RWTH Aachen University}
        {\url{https://www.rwth-aachen.de}}
        {}
        {\begin{itemize}
\item Schwerpunkt auf statistischem Lernen und algorithmischer Datenanalyse.
\item Masterarbeit zur effizienten Inferenz in tiefen neuronalen Netzen.
\end{itemize}
\textbf{Kurse}: Maschinelles Lernen, Statistische Mustererkennung, Datenbanken und Datenanalyse, Algorithmen und Datenstrukturen, Computer Vision}
\section{Berufserfahrung}

\cventry{Apr. 2021–Heute}
        {Senior KI-Ingenieur}
        {Siemens KI-Labor}
        {\url{https://www.siemens.com/ai}}
        {}
        {\begin{itemize}
\item Leitung eines Teams von 6 ML-Ingenieur:innen zur Entwicklung generativer KI-Assistenten für industrielle Anwendungen.
\item Konzeption und Umsetzung einer skalierbaren LLM-Plattform mit Retrieval-Augmented Generation und strengen Governance-Vorgaben.
\item Verantwortlich für Modellauswahl, Fine-Tuning, Evaluation und kontinuierliches Monitoring in der Cloud.
\item Enge Zusammenarbeit mit Produktmanagement und Compliance zur Einhaltung des EU AI Act.
\end{itemize}
\textbf{Schlüsselwörter}: Generative KI, LLM, RAG, MLOps, Teamführung}

\cventry{Okt. 2018–März 2021}
        {Ingenieur für Maschinelles Lernen}
        {Bosch Center for Artificial Intelligence}
        {\url{https://www.bosch-ai.com}}
        {}
        {\begin{itemize}
\item Entwicklung und Deployment von Deep-Learning-Modellen für Qualitätskontrolle und Predictive Maintenance.
\item Aufbau einer unternehmensweiten Feature-Store-Infrastruktur zur reproduzierbaren Modellierung.
\item Automatisierung von Trainingspipelines mit Kubernetes, Kubeflow und MLflow.
\item Durchführung von A/B-Tests und Modellüberwachung im Produktivbetrieb.
\end{itemize}
\textbf{Schlüsselwörter}: Deep Learning, Predictive Maintenance, Kubeflow, Feature Store, Computer Vision}

\cventry{Okt. 2016–Sept. 2018}
        {KI-Ingenieur}
        {Zalando SE}
        {\url{https://www.zalando.de}}
        {}
        {\begin{itemize}
\item Entwicklung von Empfehlungssystemen zur Personalisierung des Online-Shops.
\item Implementierung von NLP-Pipelines zur automatischen Produktklassifikation.
\item Optimierung von Modellen für Latenz und Skalierbarkeit im Hochlastbetrieb.
\item Zusammenarbeit mit Data Engineering und Product Analytics.
\end{itemize}
\textbf{Schlüsselwörter}: Recommender Systems, NLP, Python, Spark, Personalisierung}
\section{Sprachen}

\cvline{Deutsch}{Muttersprache oder zweisprachig \hfill \textbf{Schlüsselwörter}: Muttersprache}
\cvline{Englisch}{Verhandlungssichere Kenntnisse \hfill \textbf{Schlüsselwörter}: C2, Verhandlungssicher}
\cvline{Französisch}{Gute Kenntnisse \hfill \textbf{Schlüsselwörter}: B1}
\section{Fähigkeiten}

\cvline{Maschinelles Lernen}{Experte \hfill \textbf{Schlüsselwörter}: Deep Learning, NLP, Computer Vision, Generative KI, Reinforcement Learning}
\cvline{Programmierung}{Experte \hfill \textbf{Schlüsselwörter}: Python, PyTorch, TensorFlow, SQL, Bash}
\cvline{MLOps \& Cloud}{Erfahren \hfill \textbf{Schlüsselwörter}: Kubernetes, Docker, MLflow, Kubeflow, AWS SageMaker, Azure ML}
\cvline{Datenverarbeitung}{Erfahren \hfill \textbf{Schlüsselwörter}: Spark, Pandas, Airflow, dbt, Kafka}
\cvline{Führung \& Kommunikation}{Erfahren \hfill \textbf{Schlüsselwörter}: Teamführung, Agile Methoden, Stakeholder-Management, Technische Dokumentation, Mentoring}
\section{Auszeichnungen}

\cventry{Dez. 2023}
        {NeurIPS 2023}
        {Preis für den besten Beitrag}
        {}
        {}
        {\begin{itemize}
\item Auszeichnung für die Publikation „Effiziente Selbstüberwachung für domänenspezifische Sprachmodelle“.
\item Anerkennung für herausragende Beiträge zu ressourceneffizientem Training großer Sprachmodelle.
\end{itemize}}

\cventry{Juni 2022}
        {Siemens AG}
        {Innovationspreis KI}
        {}
        {}
        {\begin{itemize}
\item Preis für die Entwicklung einer skalierbaren RAG-Plattform mit messbarer Reduktion von Halluzinationen.
\item Würdigung der interdisziplinären Zusammenarbeit und des produktiven Einsatzes im industriellen Umfeld.
\end{itemize}}
\section{Zertifikate}

\cventry{März 2022}
        {Amazon Web Services}
        {AWS Certified Machine Learning – Specialty}
        {\url{https://aws.amazon.com/certification/certified-machine-learning-specialty/}}
        {}
        {}

\cventry{Mai 2023}
        {Google Cloud}
        {Google Cloud Professional Machine Learning Engineer}
        {\url{https://cloud.google.com/learn/certification/machine-learning-engineer}}
        {}
        {}

\cventry{Nov. 2020}
        {Google}
        {TensorFlow Developer Certificate}
        {\url{https://www.tensorflow.org/certificate}}
        {}
        {}
\section{Veröffentlichungen}

\cventry{Dez. 2023}
        {Effiziente Selbstüberwachung für domänenspezifische Sprachmodelle}
        {NeurIPS 2023}
        {\url{https://proceedings.neurips.cc/paper/2023}}
        {}
        {\begin{itemize}
\item Vorstellung eines selbstüberwachten Verfahrens, das ohne große Mengen annotierter Daten wettbewerbsfähige Ergebnisse erzielt.
\item Analyse der Übertragbarkeit auf ressourcenarme Sprachen und Fachdomänen.
\end{itemize}}

\cventry{Sept. 2022}
        {Skalierbare Inferenz für Generative Modelle in der Industrie}
        {KI-Fachzeitschrift}
        {\url{https://www.ki-fachzeitschrift.de/artikel/2022-09}}
        {}
        {\begin{itemize}
\item Diskussion architektonischer Muster für latenzarme Inferenz großer Sprachmodelle.
\item Praxisbericht zur Integration von RAG und Vektordatenbanken in bestehende Unternehmenslandschaften.
\end{itemize}}
\section{Projekte}

\cventry{Jan. 2023–Okt. 2023}
        {RAG-basierter Assistent zur semantischen Suche und Zusammenfassung technischer Dokumente.}
        {DocuMind – Intelligenter Dokumentenassistent}
        {\url{https://github.com/lena-hartmann/documind}}
        {}
        {\begin{itemize}
\item Entwicklung einer Ende-zu-Ende-Plattform mit Ingestion, Chunking, Embedding und Vektordatenbank.
\item Evaluation von LLMs hinsichtlich Genauigkeit, Latenz und Kosten.
\item Bereitstellung einer REST-API und eines Web-Frontends für interne Nutzer:innen.
\end{itemize}
\textbf{Schlüsselwörter}: RAG, LLM, Vektordatenbank, FastAPI, Evaluation}

\cventry{März 2019–Aug. 2020}
        {Deep-Learning-System zur automatischen Fehlererkennung in der Fertigung.}
        {VisionQC – Qualitätskontrolle mit Computer Vision}
        {\url{https://github.com/lena-hartmann/visionqc}}
        {}
        {\begin{itemize}
\item Training von Convolutional Neural Networks auf hochauflösenden Industriebildern.
\item Implementierung einer aktiven Lernschleife zur Reduktion des Annotationsaufwands.
\item Deployment auf Edge-Geräten mit TensorRT und Überwachung der Modell-Drift.
\end{itemize}
\textbf{Schlüsselwörter}: Computer Vision, Edge AI, TensorRT, Active Learning, Fertigung}
\section{Interessen}

\cvline{Open Source}{PyTorch, Hugging Face, MLflow}
\cvline{Wissenschaft \& Forschung}{NeurIPS, ICML, ACL}
\cvline{Freizeit}{Wandern, Klettern, Fotografie, Lesen}
\section{Ehrenamtliche Tätigkeiten}

\cventry{Sept. 2021–Heute}
        {Mentorin für Maschinelles Lernen}
        {AI Ethics Initiative Berlin}
        {\url{https://www.ai-ethics-berlin.de}}
        {}
        {\begin{itemize}
\item Mentoring von Nachwuchswissenschaftler:innen und Berufseinsteiger:innen im Bereich KI.
\item Organisation von Workshops zu verantwortungsvoller KI und reproduzierbarer Forschung.
\item Beitrag zur Entwicklung eines Leitfadens für ethische Evaluierung von Sprachmodellen.
\end{itemize}}
    
\end{document}